 \section{Introduction}
\label{sec:intro}

% Four moves, roughly one paragraph each. Do not exceed ~900 words for the
% capstone version; journals will want ~700.

% =====================================================================
% REVIEW 2026-09-01 . Work through these, then delete the block.

% [R4] PROTOCOL DRIFT, worth resolving before you write Move 4. The protocol
%      was git-tagged locked on 18 Aug. D26-D28 expand scope after that. The
%      forecasting work is legitimate and predeclared in the Decision Log,
%      but the locked protocol and the manuscript currently describe
%      different studies. Add an amendment-log entry in protocol section 19
%      so the sequence is on the record.
% [R5] No citations in the drafted prose. The scaffold asks for the range
%      assessment and the global decline synthesis. Nothing in the
%      Literature Log covers the bottleneck dating or the range-decline
%      claim yet, so those need sourcing before they can stay.
% =====================================================================

The southern African cheetah has had their range substantially reduced, with their current range estimated to be on majority unprotected land. Cheetah range on protected land is governed by centralized park-management initiatives \parencite{buk2018metapopulation}. Unprotected land poses a different type of challenge, with landowners, local communities and governmental entities involved. This paper looks at landscape connectivity change between 2012 and 2024. The core range of cheetahs is fixed at the 2010-2016 baseline, so modeled changes in permeability are attributable to landscape changes rather than changes in the core range itself.

\subsection*{Move 2: What is already modelled, and what is not}

Previous research in the region has predominantly focused on habitat suitability, a topic that is already comprehensively addressed in the existing literature. The objective of my study is not to develop an additional suitability model, but rather to identify and evaluate candidate spatial “pinch points” (i.e., constricted linkages) among core areas. The study most comparable to the present work is a recent multi-species connectivity assessment; however, that analysis explicitly excluded cheetah due to insufficient data. Data limitations were therefore explicitly considered and mitigated to the extent possible in this study, and they provide both the primary justification for the existence of this paper and its principal constraint.

\subsection*{Move 3: The methodological frame}

Two primary connectivity approaches are commonly used to model these kinds of frameworks: least-cost connectivity and circuit-theoretic connectivity (\cite{adriaensen2003leastcost}; \cite{mcrae2008circuit}). The least-cost approach links core areas by selecting the route with the lowest cumulative “cost,” meaning the path of least resistance between cores, based on weight values assigned to each multivariate resistor. In contrast, the circuit-theoretic approach represents the landscape as a resistor network, sends a “current” through it, and highlights redundancy across the broader landscape, relying on the same multivariate resistors. A key limitation of both methods is that you must first assign a resistance “score” or “weight” to every cell (1 km in my case). \textcite{zeller2012resistance} demonstrated that even reasonable alternative assumptions about resistance values can influence results more strongly than the choice of connectivity method itself. For that reason, rather than treating one resistance surface as definitive, I varied the weights along two axes. The first changed the balance among the anthropogenic pressures, comparing my primary weighting against an equal weighting and an infrastructure-emphasis weighting. The second changed how much the surface leaned on vegetation relative to anthropogenic pressure and terrain, at 10, 20 and 40 percent. The vegetation-balanced surface combined with the documented fence treatment carries the primary results, and the remaining combinations are reported as sensitivity rather than as alternative answers. A link is only described as robust here if it survives them. This helps to minimize bias from the selection of resistor weights.

\TODO{Sentence above now matches what was actually run. Two small things left:
1. The protocol still calls these RES-01 to RES-06. Worth a line in the protocol saying
   the implemented set differs, mostly so your own record stays straight.
2. Cites worth adding here: \textcite{zeller2018equal} alongside
   \textcite{zeller2012resistance}, and \textcite{beier2008forks} for the
   design-choices point.}
\subsection*{Move 4: Objectives and hypotheses}
% Sentence starters: "The first objective was to ..." / "I predicted that ..." / "This prediction would be unsupported if ..."

The first objective was to measure changes in modeled structural connectivity among fixed-range cores across four time points from 2012--2024. The second was to determine candidate pinch points and linkages that remain influential across the tested resistance assumptions and describe their protected-area status. I made the following not mutually exclusive predictions in support of each hypothesis.



\textbf{H1 (revised from the workbook).} Connectivity loss is concentrated on a
minority of network links rather than distributed across the network, and the
affected links are disproportionately those with high betweenness in the
baseline network.

\emph{Note on the revision.} H1 as originally drafted predicted that cells with
larger increases in human pressure would show larger increases in resistance.
That is true by construction, since the pressure layers are the resistance
inputs. It cannot fail and therefore is not a hypothesis. The revised form asks
a question the model can answer negatively: pressure change might fall on
redundant links, in which case network-level connectivity is robust to it and
the conservation implication reverses. 

\emph{Refuted if:} connectivity loss is approximately uniform across links, or
loss falls preferentially on low-betweenness links.

\textbf{H2.} High-value connectivity area is disproportionately located outside
formal protected areas.

\emph{Refuted if:} the outside-PA share of high-current area does not exceed the
predeclared null (\S\ref{sec:pa}) at the stated significance level, or the
result does not hold across current-flow thresholds.

\textbf{H3.} A limited subset of transboundary linkages remains important across
all resistance assumptions, while other apparent corridors are strongly
model-sensitive.

\emph{Refuted if:} agreement is uniformly high (in which case the ensemble adds
nothing) or uniformly low (in which case no robust set exists and the paper's
primary output becomes the uncertainty map, per the stopping rule in the
predeclared analysis protocol).

\paragraph{Terminology.} Throughout, \structconn\ refers to modelled landscape
permeability between fixed range cores. It is not evidence of cheetah movement,
dispersal, or gene flow. Where a location is described as a pinch point, this means
a location of concentrated modelled current or a link-level bottleneck, not an
observed crossing.
% This paragraph is CAU-001's mitigation. Keep it in the Introduction, not just
% the Discussion. It cannot be cut for length.
% REVIEW: agreed, and it is the clearest statement of your claim boundary
% anywhere in the manuscript. It must survive every length cut you make.

% =====================================================================
% REVIEW 2026-09-01 . CRAFT PATTERN (the thing to carry into Move 2-4)
%
% Your costliest habit is choosing verbs that claim more than the method
% delivers: "observed" change, cheetahs that "will be roaming", a
% "predictive look" at permeability. Each is a small word doing large
% unearned work. In a paper whose entire contribution rests on separating
% modelled permeability from real movement, these are not style slips,
% they are the failure mode itself.
%
% THE RULE, apply it as you write rather than in revision:
% before every verb describing what you did or what happens, ask whether a
% sensor or a person recorded it, or whether your model produced it. If the
% model produced it, the verb is "modelled", "estimated" or "inferred", and
% the noun takes "modelled" as a modifier. Reserve "observed" for the Dryad
% records and the Limpopo tracking data, which are the only observations in
% this project.
%
% Apply it to your TODO notes too. The scaffold currently contains the same
% error ("observed change in network permeability"), so it is teaching you
% the habit every time you read it.
% =====================================================================
